\subsection{Conclusão}
O TF-IDF considera a informação de toda a resenha e obteve uma acurácia de 90,14\%, enquanto a TextCNN (que se limitou no nosso estudo a 400 tokens) alcançou acurácia de 86,78\%. Ao utilizar o texto completo, mantendo os mesmos pesos da CNN, observamos uma melhora de 0,50 ponto percentual na acurácia geral e de 4,66 pontos percentuais nas resenhas com mais de 500 palavras. No entanto, a diferença de desempenho entre os modelos também aparece nos textos curtos. 
Dessa forma, concluímos que o tamanho da resenha para a TextCNN (quantidade de tokens considerados) é uma oportunidade a ser explorada de acordo com o poder computacional e tempo disponível, assim como a arquitetura da rede e todos os demais hiperparâmetros. Em todo o caso, o modelo baseline mostrou um ótimo desempenho para o problema abordado, principalmente tendo em vista sua baixa complexidade em comparação com um modelo de deep learning.  

Além disso, o experimento com janelas convolucionais iguais mostrou que o Global Max Pooling perde informações sobre a posição dos padrões identificados. Já no exemplo com \textit{but}, a mudança na ordem das palavras altera as janelas analisadas, e a classificação depende tanto das ativações dos filtros quanto da camada final.

Por fim, o BERT ajustado (90,94\%) e o Qwen com pesos congelados (95,26\%) superaram o baseline, embora utilizem arquiteturas, representações de entrada e pré-treinamentos diferentes.
\fonteatual{seminario/revisao_metodologica.md}{Controles}.

\begin{thebibliography}{9}
\footnotesize
\setlength{\itemsep}{0pt}
\setlength{\parskip}{0pt}
\bibitem{maas} Maas, A. L. et al. \textit{Learning Word Vectors for Sentiment Analysis}.
ACL, 2011. \href{https://aclanthology.org/P11-1015/}{ACL Anthology}.
\bibitem{kim} Kim, Y. \textit{Convolutional Neural Networks for Sentence Classification}.
EMNLP, 2014. \href{https://aclanthology.org/D14-1181/}{ACL Anthology}.
\bibitem{sennrich} Sennrich, R., Haddow, B. e Birch, A.
\textit{Neural Machine Translation of Rare Words with Subword Units}. ACL, 2016.
\href{https://aclanthology.org/P16-1162/}{ACL Anthology}.
\bibitem{prince} Prince, S. J. D. \textit{Understanding Deep Learning}.
MIT Press, 2023. Caps. 6, 8, 9 e 10. \href{https://udlbook.github.io/udlbook/}{Livro e materiais}.
\bibitem{devlin} Devlin, J. et al. \textit{BERT: Pre-training of Deep Bidirectional
Transformers for Language Understanding}. NAACL, 2019. \href{https://aclanthology.org/N19-1423/}{ACL Anthology}.
\bibitem{gardner} Gardner, M. et al. \textit{Evaluating Models' Local Decision
Boundaries via Contrast Sets}. Findings of EMNLP, 2020.
\href{https://aclanthology.org/2020.findings-emnlp.117/}{ACL Anthology}.
\bibitem{kotzias} Kotzias, D. et al. \textit{From Group to Individual Labels Using Deep Features}.
KDD, 2015. \href{https://doi.org/10.1145/2783258.2783380}{Artigo}.
\href{https://doi.org/10.24432/C57604}{Coleção UCI}.
\bibitem{qwen} Yang, A. et al. \textit{Qwen3 Technical Report}, 2025.
\href{https://arxiv.org/abs/2505.09388}{arXiv}.
\end{thebibliography}
